\section{Loss of second derivatives at positive exponent}
\label{sec:numa-positive-exponent-regularity}

Let $X=\mathbb P(\mathbb C^2\oplus\mathbb C^2)$, with block coordinates $[A:B]$, and let $\omega_X$ be the Fubini--Study form. On $U=\{A\ne0,\ B\ne0\}$ use the radial coordinate $r$ and rank-one forms $P,Q,R$ of \cref{subsec:numa-rank-directions}. Take the scalar coefficients $p,q,b,z_0,x,y,z$ from \cref{subsec:numa-scalar-coefficients}. Let $\omega_0=xP+yQ+zR$, and let $\rho$ and $\omega=\omega_0+\ddc\rho$ be the correction and Hermitian form from \cref{subsec:numa-positive-background}. The density $f$ is the common cube there, so $f\omega_X^3=0$ on $|r|\leq1$ and $f\omega_X^3=\omega_0^3$ outside that band. Write $\mathrm{MA}_\omega(u)=(\omega+\ddc u)^3$ for bounded $\omega$-plurisubharmonic $u$. Choose the correction cutoff even and let $K$ be the group generated by block unitary transformations and block exchange. Let $t_*\in(0,1)$ be the unique solution of $\int_0^{t_*}z_0(s)\,\mathrm{d}s=p(t_*)$, and let $u_*$ be the normalized symmetric solution of \cref{prop:numa-radial-continuum}. In the flux formula below use $A(r,s)=(x(r)+s)(y(r)-s)$ and $k(r)=0$ on $|r|\leq1$, $k(r)=x(r)y(r)z(r)$ on $|r|>1$, as in \cref{lem:numa-radial-flux}.

The selected solution is not $C^2$. In fact the same local loss of regularity occurs at every positive exponent, despite smoothness of the background and the nonnegative density. Thus the exponential term restores uniqueness without removing this degeneracy.

\Needspace{12\baselineskip}
\begin{proposition}\label{prop:numa-positive-exponent-nonc2}
For the symmetric projective data of \cref{thm:numa-selected-limit}, let $u_\lambda$ solve
\[
 \mathrm{MA}_\omega(u_\lambda)
       =\mathrm{e}^{\lambda u_\lambda}f\omega_X^3,
 \qquad\lambda>0.
\]
On the open band $\{-1<r<1\}$, $u_\lambda-u_*$ is constant. Both potentials are $C^{1,1}$ there and fail to be $C^2$ at every point of the two hypersurfaces $r=\pm t_*$. In particular, none of these positive-exponent solutions is a globally $C^2$ potential.
\end{proposition}
\begin{proof}
Uniqueness and invariance of the data imply that $u_\lambda$ is $K$-invariant. Write $u_\lambda+\rho=G_\lambda(r)$ on $U$. The positivity argument in \cref{lem:numa-radial-flux} uses only boundedness and invariance, so $G_\lambda$ is locally Lipschitz and $h_\lambda=G_\lambda'$ is locally of bounded variation, with
\[
 -x\leq h_\lambda(r\pm)\leq y,
 \qquad \mathrm{d}h_\lambda+z\,\mathrm{d}r\geq0.
\]
Its weak flux equation is the same one, with the prescribed coefficient replaced by
\[
 k_\lambda(r)=\mathrm{e}^{\lambda(G_\lambda(r)-\rho(r))}k(r).
\]
Indeed, the scalar convolution and signed-product argument of that lemma is unchanged, and this is exactly the radial coefficient of the right-hand measure. At the initial bounded stage $k_\lambda$ is locally bounded. Hence the right side of the flux equation has no atoms. The same strictly positive jump integral eliminates every jump of $h_\lambda$. Thus $h_\lambda$ is continuous and $G_\lambda$ is $C^1$. The coefficient $k_\lambda$ is consequently continuous. Where $A(r,h_\lambda)>0$, inversion of the $C^1$ flux gives
\[
 h_\lambda'=\frac{k_\lambda(r)}{A(r,h_\lambda)}-z.
\]

On $-1<r<1$ the coefficient $k_\lambda$ is identically zero. Exchange invariance gives $h_\lambda(0)=0$ and oddness. On the positive half-band, the first-contact and no-departure arguments in the proof of \cref{thm:numa-selected-limit} apply verbatim: the first branch solves $h'=-z_0$, reaches $-p$ at $t_*$, and cannot leave that lower boundary to its right because $(h+p)'=p'-z_0<0$ on every proposed departure interval. Therefore
\[
 h_\lambda(r)=
 \begin{cases}
  p(r),&-1<r\leq-t_*,\\
  -\displaystyle\int_0^r z_0(s)\,\mathrm{d}s,&-t_*\leq r\leq t_*,\\
  -p(r),&t_*\leq r<1.
 \end{cases}
\]
This is the derivative of $u_*+\rho$. Its primitive differs from $G_\lambda$ by one constant on the connected band. The claimed local $C^{1,1}$ regularity follows from the explicit slope.

At $t_*\in(0,1)$ the two one-sided derivatives of this slope are
\[
 h_\lambda'(t_*-)=-z_0(t_*)<0,
 \qquad
 h_\lambda'(t_*+)=-p'(t_*)>0.
\]
These unequal derivatives prevent the scalar primitive from having a second derivative at $t_*$. Through every point of that level set, the submersion $r$ supplies a smooth curve parametrized by $r$, transferring the failure to $u_\lambda$. Smooth $\rho$ cannot cancel the jump. Oddness gives the conclusion at $-t_*$, and the identical slope proves it for $u_*$.
\end{proof}

The zero set of $f$ contains the whole band. This example therefore does not concern the smooth strictly positive densities in the Hermitian Calabi--Yau theorem. It does show why smoothness of a nonnegative density alone is insufficient for second-order regularity, even when the equation includes a strictly increasing exponential factor.
